%! AP Statistics — Unit 7, Lesson 7.3 — Group Activity
\documentclass[10pt]{article}
\usepackage{apstats-article}
\usepackage{apstats-boxes}

\begin{document}

\pageheader{Unit 7, Lesson 7.3}{Group Activity}
\namepartnerperiod

\vspace{0.1in}
\textbf{Directions:} Three completed confidence intervals are given below. For each,
complete all parts.

\begin{tcolorbox}[colback=white, colframe=black!40, title=\textbf{Tier R --- Remediate},
    fonttitle=\bfseries, arc=2mm, left=3mm, right=3mm, top=2mm, bottom=2mm]

\textbf{CI 1:} A random sample of 30 eighth graders yields a 95\% $t$-interval for
mean math test score: $\mathbf{(72.1,\ 84.3)}$ points.

\begin{enumerate}[label=\alph*.]
  \item Write a correct interpretation. Use the sentence frame:\\[4pt]
        ``We are \blank{2cm}\% confident that the interval (\blank{2cm}, \blank{2cm})
        \blank{1cm} captures the true mean \blank{5cm} for \blank{5cm}.''

  \item The district claims the mean score is 70 points. Does the interval provide
        convincing evidence against this claim? Circle: \quad \textbf{Yes} \quad \textbf{No}

        \vspace{0.05in}
        Explanation (one sentence): \writelines{2}

  \item If the sample size were 120 instead of 30, would the margin of error be
        larger, smaller, or the same? By approximately what factor?

        \writelines{2}
\end{enumerate}
\end{tcolorbox}

\begin{tcolorbox}[colback=white, colframe=black!40, title=\textbf{Tier A --- Approaching Proficiency},
    fonttitle=\bfseries, arc=2mm, left=3mm, right=3mm, top=2mm, bottom=2mm]

\textbf{CI 2:} A random sample of 25 adult runners yields a 90\% $t$-interval for
mean 5K finish time: $\mathbf{(24.6,\ 29.8)}$ minutes.

\begin{enumerate}[label=\alph*.]
  \item Write a correct interpretation. \writelines{2}

  \item A coach claims the mean finish time for all adult runners is 30 minutes.
        Does the interval provide convincing evidence against this claim? Justify
        using the interval. \writelines{2}

  \item The researcher wants to raise the confidence level to 99\% for a follow-up
        study (same sample). Describe the effect on the width of the interval and
        the margin of error. \writelines{2}

  \item A different researcher claims the mean is greater than 25 minutes. Does the
        interval support this claim? Explain. \writelines{2}
\end{enumerate}
\end{tcolorbox}

\begin{tcolorbox}[colback=white, colframe=black!40, title=\textbf{Tier E --- Extension},
    fonttitle=\bfseries, arc=2mm, left=3mm, right=3mm, top=2mm, bottom=2mm]

\textbf{CI 3:} A product manager constructs a 95\% $t$-interval for the mean battery
life of a new phone model: $\mathbf{(18.4,\ 21.6)}$ hours ($n = 50$, $s = 5.8$ hours).

\begin{enumerate}[label=\alph*.]
  \item Write a correct interpretation. \writelines{2}

  \item The company advertises ``at least 20 hours of battery life on average.''
        Does the interval provide convincing evidence against this advertising claim?
        Justify carefully. \writelines{2}

  \item The product manager wants the margin of error to be at most 1 hour. What
        minimum sample size would be needed? (Use $s = 5.8$ and $t^* \approx 1.960$
        as a starting estimate.) Show your work. \writelines{4}

  \item Describe the trade-off if the manager instead raises the confidence level
        from 95\% to 99\% (keeping $n = 50$). \writelines{2}
\end{enumerate}
\end{tcolorbox}

\end{document}
